\clearpage
\fancyhead[L]{\small PIAZZI. LEZIONI DI ASTRONOMIA.}
\hypertarget{native-622-lower}{}
\begin{center}\small Göttingische gelehrte Anzeigen. Stück 4. Seite 33..37. 1820 Januar 6.\\[3pt]\rule{.72\linewidth}{.3pt}\end{center}

\noindent\emph{Lezioni elementari di astronomia ad uso del real osservatorio di Palermo. Palermo {\upshape 1817}. Dalla stamperia reale.} Tomo I, 240 S.; Tomo II, 446 S. in kl. 4.

Der Verf. dieses Lehrbuchs, der berühmte und verdiente \textsc{Piazzi}, erklärt in der Vorrede, dass dasselbe zunächst auf Veranlassung der Vorlesungen, welche von ihm auf der Palermer Sternwarte gehalten worden, entstanden und eigentlich nicht für das Publicum bestimmt gewesen sei, und dass auf den dringenden Wunsch der Schüler die Vorgesetzten des öffentlichen Unterrichts die Bekanntmachung befohlen haben. Von Schriften dieser Art erwartet man keine eigentliche Erweiterung der Wissenschaft selbst; eine klare, gründliche, der Bestimmung angemessene Anordnung des Vortrags ist die Hauptsache. Die Bedürfnisse des Vortrags sind wieder nach Massgabe der Vorkenntnisse und Absichten der Schüler sehr verschieden; anders für blosse Dilettanten, anders für mathematisch gebildete; anders für den, welcher Sinn dafür hat, in der Wissenschaft eines der vollkommensten und edelsten Producte des menschlichen \emph{Geistes} zu verehren, der Nichts auf Treue und Glauben annimmt, sondern überall streng logischen Zusammenhang verlangt, und dem die vollendetste Darstellung die liebste ist, wenn sie auch mehr als gewöhnliche Anstrengung und Hülfskenntnisse erfordert, und endlich wieder anders für den, welcher nur nach dem niedern Ziele strebt, sich zur Beobachtung und Berechnung himmlischer Phänomene, so viel es eine beschränktere mathematische Vorbildung gestattet, geschickt zu machen. Der Verf. gibt zwar nicht ausdrücklich an, für welche Classe von Schülern er dieses Buch eigentlich bestimmt; allein es scheint doch, dass er sich vornehmlich solche, die zur letzten Classe gehören, gedacht habe. Von einem Astronomen, wie \textsc{Piazzi}, der um mehrere Theile der Astronomie so grosse Verdienste hat, liess sich übrigens erwarten, dass auch ein elementarisches Werk dieser Art eine Menge eigenthümlicher Ansichten, merkwürdiger Aeusserungen und Urtheile enthalten musste, und in dieser Beziehung ist es auch selbst für Astronomen nicht ohne Interesse.

Der Verf. hat sein Werk in sieben Bücher abgetheilt, wovon das erste eine allgemeine Uebersicht über die himmlischen Erscheinungen gibt, so weit man dazu blos durch rohe Beobachtung, ohne Instrumente und ohne besondere mathematische Kenntnisse gelangen kann. Es wird hier also gehandelt von der allgemeinen täglichen Bewegung, von den auf der Himmelskugel angenommenen Punkten und Kreisen, von der Gestalt der Erde, von dem Unterschiede zwischen Fixsternen und Planeten, von der Sonne, dem Monde, den Planeten und Cometen, von der Ekliptik und den übrigen dadurch bestimmten Kreisen der Himmelskugel, vom Thierkreise und den Sternbildern, von der Vorrückung der Nachtgleichen, von den Finsternissen, von den ersten Ungleichheiten der Planeten und ihrer Erklärung durch die Kopernicanische Weltordnung. Das zweite Buch beschäftigt sich mit Vorkenntnissen der neuen Astronomie. Wir finden hier etwas über die Auflösung sphärischer Dreiecke, eine Uebersicht der vornehmsten, heutzutage gebräuchlichen astronomischen Werkzeuge, als Fernrohre, Loth, Libelle, Verniers, Micrometer, Mittagsternrohr, Uhr, Vollkreis, Quadrant, Zenithsector, Aequatoreal und Vervielfältigungskreise (über letztere urtheilt \textsc{Piazzi} nicht günstig; indessen scheint er doch das, was zu ihrem \SourceLeaf{623}Nachtheile gesagt werden kann und auch von andern gesagt ist, etwas übertrieben zu haben. Hätte er selbst, sagt \textsc{Piazzi}, mit einem Instrument dieser Art beobachten sollen, so würde er \emph{nicht den hundertsten Theil} dessen, was er gethan, geleistet haben). Hiernächst wird von der Refraction, der Parallaxe, der Aberration und Nutation gehandelt, in so fern rücksichtlich dieser Umstände die Beobachtungen einer Verbesserung bedürfen; dann von der Zeit, und den Methoden der Zeitbestimmung; endlich von den Beobachtungen, die zur Bestimmung der Polhöhe, der Schiefe der Ekliptik und der Stellung der Gestirne gegen Aequator und Ekliptik dienen. Das dritte Buch beschäftigt sich dann ausschliesslich mit den Fixsternen. Sehr ausführlich über die Berechnung der Präcession, Nutation und Aberration, dann über die jährliche Parallaxe der Fixsterne, ihre eigene Bewegung, über die Sternverzeichnisse und die eigenthümlichen Merkwürdigkeiten der Fixsterne. Einen sonderbaren Misgriff bei Bestimmung der Grösse der jährlichen Präcession können wir hier nicht unbemerkt lassen. \textsc{Piazzi} zieht von der beobachteten jährlichen Aenderung der Länge von 15 Fixsternen, die aus der Vergleichung von \textsc{Bradley}'s Bestimmung für 1755 mit \textsc{Piazzi}'s eigner für 1805 folgt, theils die Wirkung der Verrückung der Ekliptik, theils die der \emph{eignen Bewegung} ab, und nimmt dann aus allen (übrigens sehr ungleich ausfallenden) Resultaten das Mittel. Dies ist ein offenbarer logischer Cirkel, indem es kein anderes Mittel gibt, die eigne Bewegung der einzelnen Sterne zu erkennen, als wenn man die beobachtete Bewegung mit der schon als bekannt angenommenen Präcession vergleicht. Bei einer folgerechten Rechnung hätten alle Sterne einerlei Resultat, und zwar dasselbe wieder geben sollen, wovon man ausgegangen war, um die eigenen Bewegungen zu finden.

Der Gegenstand des vierten Buches ist die elliptische Bewegung der Planeten, grösstentheils aus der Theoria motus corporum coelestium entlehnt, jedoch mit Weglassung aller Beweise.

Im 5ten Buche werden die Himmelskörper unsers Sonnensystems einzeln betrachtet. Merkwürdig ist die Art, wie \textsc{Piazzi} sich S. 182 über die Bestimmung der Rotationszeit der Venus durch den verstorbenen \textsc{Schröter} äussert. `Niente di meno non manca, chi non poco dubiti di tali scoperte, e le consideri come non molto dissimili dall' altra del satellite, che \textsc{Short} e qualche altro credeva di aver veduto intorno a quaesta pianeta. Una illusione ottica, un fenomeno non benè osservato, accompagnati da un poco di immaginazione, hanno talora fatto avanzare delle congetture, e sognare delle esseri, che ben presto son poi caduti in piena dimenticanza (\textsc{Piazzi} scheint hier doch vergessen zu haben, dass er selbst S. 46 des ersten Bandes das Dasein eines Venus-Trabanten für sehr wahrscheinlich erklärt hatte, an den jetzt schwerlich irgend ein anderer Astronom noch glaubt). Eben so erklärt er sich S. 230 gegen die von einigen Astronomen angeblich gesehenen brennenden Mondvulkane.

Im sechsten Buche wird von den Finsternissen, Sternbedeckungen und Vorübergängen der untern Planeten vor der Sonnenscheibe gehandelt, und endlich im siebenten von den Cometen, zu deren Bahnbestimmung hier nur die indirecte Methode in derselben höchst unvollkommenen und man kann sagen, planlosen Gestalt gegeben wird, die die Astronomen zu ihrer Verwunderung in einem bekannten, vor einigen Jahren erschienenen Lehrbuche der Astronomie gefunden haben.
